\section{Problem Formulation}
\label{sec:model}

\subsection{Applications, roles, and observations}

We model a web application under test (AUT), as seen by a fixed role $r$, as a
labelled transition system $M_r=(Q,A,\delta,q_0)$ over concrete UI
configurations $Q$, actions $A$, and a partial transition function
$\delta:Q\times A\rightharpoonup Q$. The auditor never sees $Q$ directly. A
\emph{capture} $\kappa(q)=(u,d,t,s)$ bundles the URL, serialized DOM,
accessibility (AX) tree, and an optional screenshot. \sys{} pins one role per
run; multi-role coverage is several runs, and comparing roles is a
reporting-layer concern (\cref{sec:discussion}).

\begin{definition}[State abstraction]
A \emph{fingerprint} is $\phi = H\circ N$, where $N$ normalizes a capture by
removing volatile content (session and CSRF tokens, nonces, long numeric
timestamps, generated hex identifiers, per-session CDP node ids, transient
focus), and $H$ is SHA-256. Two configurations are \emph{the same abstract
state} iff $\phi(\kappa(q))=\phi(\kappa(q'))$.
\end{definition}

Normalization has two failure modes. If it is too sensitive, the model
explodes into one node per page load. If it is too coarse, a security-relevant
change (for example, a destructive button that appears) folds into an existing
state and is never audited. \Cref{sec:eval:fp} measures both.

\subsection{The exploration graph is cyclic}

\begin{definition}[Exploration graph]
The \emph{exploration graph} is a directed multigraph $G=(V,E)$ with
$V\subseteq\operatorname{img}\phi$ and $E\subseteq V\times A\times V$, where
$(v,a,v')\in E$ records that executing $a$ in $v$ was observed to lead to $v'$.
An edge is a \emph{back-edge} if $v'$ was already discovered when it was added.
\end{definition}

Model-inference crawlers often treat revisits as redundant and drop them,
leaving a tree or DAG. We treat back-edges as \emph{evidence}. ``Log out
returns to login'' and ``Admin returns to Dashboard'' are facts about the
AUT's navigation structure, and a path-sensitive audit can only reason over
them if they are kept. \sys{} therefore never prunes an edge. Loop prevention
lives entirely in a \emph{claim set} $C\subseteq V\times A$.

\begin{definition}[Claim-before-execute]
Action $a$ is executed from $v$ only if $(v,a)\notin C$, and $(v,a)$ is inserted
into $C$ \emph{before} $a$ is sent to the browser. The edge $(v,a,v')$ is
persisted only \emph{after} $a$ succeeds.
\end{definition}

\begin{proposition}[Termination and at-most-once per pair]
\label{prop:term}
Suppose $|V|\le n$, which is enforced by a state cap, and each state enables
at most $k$ actions. Then exploration performs at most $nk$ \emph{frontier}
executions and terminates. No claimed pair $(v,a)$ is executed as a frontier
action twice, including across crash and resume.
\end{proposition}
\begin{proof}[Proof sketch]
Every frontier execution inserts a distinct pair into $C$ first, and
$|C|\le nk$. On resume, a pair that is in $C$ but has no edge is treated as
done and never re-fired. Execution strictly follows insertion, so a crash
cannot cause a second execution of the same pair.
\end{proof}

Most AUTs are not idempotent: re-submitting a form changes the application
being audited. Reaching a queued action with a single browser, however, means
\emph{replaying} its path from the seed, and replay re-fires the actions on
that path. \Cref{prop:term} is therefore a guarantee about \emph{pairs}, not
about \emph{effects}. \sys{} narrows the gap by kind: only
$\textsf{ReplaySafe}=\{\texttt{navigate},\texttt{click},\texttt{back}\}$ may be
re-fired. A candidate whose path passes through a \texttt{fill},
\texttt{select}, or \texttt{submit} is recorded as skipped, with a reason, and
its subtree is reported as unexplored instead of being reached by quietly
re-submitting a form. The residual assumption is that \emph{clicks are
navigational}. A click that mutates state on a replayed path (for example, a
one-click ``Delete'') \emph{is} re-fired. We treat this as an open limitation
(\cref{sec:discussion}), not a solved problem. Action identity is
\emph{content-addressed}:
$\mathit{id}(a)=H(\mathit{role}\,\|\,\mathit{name}\,\|\,\mathit{selector})$, not
a DOM index. Replay therefore survives DOM reordering, and the same ``Log out''
control keeps one identity everywhere, while deduplication stays scoped to
$(v,a)$.

\subsection{Policies as two set operations}

A policy is written in English by a QA engineer and compiled into
machine-checkable clauses $c=(\sigma_c,\rho_c,\pi_c)$: a \emph{subject}
$\sigma_c$ (a UI element kind or capability), a role scope $\rho_c$ (empty
means all roles), and a polarity $\pi_c\in\{\textsf{must\_exist},
\textsf{must\_not\_exist}\}$. This splits the policy into forbidden clauses
$\Forb$ and required clauses $\Req$. Perception maps each state to a set of
grounded subjects $\Sig(v)$ (element tags, accessible names, and capabilities).
For a role $r$, let $\Forb_r$ and $\Req_r$ be the clauses whose scope admits
$r$.

\begin{definition}[Negation-aware audit]
\begin{align}
\mathit{Viol}_\Forb(v) &= \{\,c\in\Forb_r : \sigma_c\in\Sig(v)\,\}
  && \text{(forbidden present)}\label{eq:forb}\\
\mathit{Viol}_\Req(v)  &= \{\,c\in\Req_r : \sigma_c\notin\Sig(v)\,\}
  && \text{(required absent)}\label{eq:req}
\end{align}
The audit of $v$ is $\mathit{Viol}_\Forb(v)\cup\mathit{Viol}_\Req(v)$, and each
violation is tagged with the operation that produced it.
\end{definition}

\Cref{eq:forb} is an intersection ($S\cap\Forb$) and \cref{eq:req} is a
difference ($\Req\setminus S$). No single intersection can express both, so an
intersection-only oracle is structurally unable to report a missing element
(\cref{sec:eval:rq3}).

\subsection{Error-direction asymmetry under imperfect perception}
\label{sec:model:asym}

In a real system, $\Sig(v)$ is replaced by a perception estimate $\hat\Sig(v)$
from a DOM/AX correlator or a vision-language model. The two clause types react
to perception errors in \emph{opposite} directions.

\begin{proposition}[Error-direction asymmetry]
\label{prop:asym}
Let $\widehat{\mathit{Viol}}$ denote the audit computed on $\hat\Sig$.
\begin{enumerate}
  \item If perception \emph{under-reports} ($\hat\Sig(v)\subseteq\Sig(v)$), then
    $\widehat{\mathit{Viol}}_\Forb(v)\subseteq\mathit{Viol}_\Forb(v)$ and
    $\widehat{\mathit{Viol}}_\Req(v)\supseteq\mathit{Viol}_\Req(v)$. Forbidden
    clauses can only be missed, and required clauses can only raise false
    alarms.
  \item If perception \emph{over-reports} ($\hat\Sig(v)\supseteq\Sig(v)$, for
    example by hallucinating a subject), both inclusions reverse.
\end{enumerate}
\end{proposition}
\begin{proof}
\Cref{eq:forb} is monotone in $\Sig$ and \cref{eq:req} is antitone in $\Sig$.
\end{proof}

\begin{corollary}
No uniform perception bias is safe for both clause types. A detector tuned to
avoid false alarms on forbidden clauses (conservative, under-reporting) creates
false alarms on required clauses, and the reverse also holds.
\end{corollary}

This matters for multimodal perception. Vision-language models answer negated
queries close to chance~\cite{alhamoud2025negation}, and their compositional
understanding is known to be shallow~\cite{yuksekgonul2023aro,ma2023crepe}.
Those studies use natural images, not UI screenshots, so we cite them as
motivation rather than as a measurement of UI perception. Whatever a
perceiver's error rate, \cref{prop:asym} shows that a single yes/no ``is
$\sigma$ on the screen?'' query has a polarity-dependent risk profile. This motivates the design
choices in \cref{sec:system:negation}: evidence-carrying verdicts, a
deterministic DOM/AX channel that cross-checks the VLM, and (proposed) a
three-valued verdict that abstains instead of guessing when the two channels
disagree. \Cref{sec:eval:rq2} observes exactly the asymmetry predicted by
\cref{prop:asym}.
